% Satz des Handbuchs: KOMA-Script, LuaLaTeX, Libertinus und Source Sans
\usepackage{fontspec}
% Zeichen, die Libertinus nicht hat (⊙ gestempelt, ▾), kommen aus DejaVu
\directlua{luaotfload.add_fallback("ersatz", {"DejaVu Sans:mode=node;"})}
\setmainfont{Libertinus Serif}[RawFeature={fallback=ersatz}]
\setsansfont{Source Sans Pro}[RawFeature={fallback=ersatz}]
\setmonofont{DejaVu Sans Mono}[Scale=0.85]
\usepackage{xcolor}
\definecolor{phila}{HTML}{7A4A1E}
\addtokomafont{disposition}{\color{phila}}
\addtokomafont{title}{\color{phila}}
\usepackage{float}
\floatplacement{figure}{H}
\usepackage[font=small,labelformat=empty]{caption}
\AtBeginDocument{\hypersetup{colorlinks=true, linkcolor=phila, urlcolor=phila}}
\setlength{\emergencystretch}{3em}
\titlehead{\centering\includegraphics[width=4cm]{icon/symbol.png}}
\publishers{\small Freund der Philatelie · Briefmarken als Zeitgeschichte\\[1ex]\url{https://philaphil.bgg-home.de/}}
\usepackage{scrlayer-scrpage}
\pagestyle{scrheadings}
\clearpairofpagestyles
\ihead{\small\color{gray}PhilaPhil · Benutzerhandbuch}
\ohead{\small\color{gray}\headmark}
\cfoot{\small\pagemark}
\automark{section}
